% ============================================================
% 65_code_style.tex — OPTIONALES Code-Listing-Modul (CodeExpert)
% ============================================================
% Opt-in: in preamble.tex einkommentieren für code-lastige
% Fächer (Informatik). Das Modul lädt seine optionale Abhängigkeit
% listings samt tcolorbox-Listings-Library selbst.
% Farben (ce_*, code_*) aus 40_colors_structure.tex,
% Längen (\ZSFcode*) aus 30_layout_spacing.tex,
% Box-Grundvertrag (zsftitlebox) und Regler aus 60_boxes.tex.
%
% Stellt bereit:
%   - lstlisting-Style "CodeExpert" (dunkler Code-Block)
%   - codebox[Titel][Optionen]  — Container für Code-Snippets; ohne
%                                 vertikale Polsterung, weil lstlisting
%                                 eigene Skips mitbringt
%   - Regler `part`             — whole (Default) | first | mid | last:
%                                 mehrteilige Code-Gruppen mit offenen
%                                 Kanten, zwischen denen die Spalte
%                                 umbrechen darf
% ============================================================

\RequirePackage{listings}
\tcbuselibrary{listings}

% -------------------------
% CodeExpert Listings Style
% -------------------------
% Was NICHT zum Stil gehört: Satzmechanik, die für jedes Listing gilt.
% Alles, was CodeExpert unten ohnehin setzt, stand hier ein zweites Mal —
% andere Rahmenart, andere Skips — und welcher Wert galt, hing daran, ob der
% Autor den Stil angefordert hatte.
\lstset{
    columns=flexible,
    tabsize=2,
}

\lstdefinestyle{CodeExpert}{
    language=Python,
    basicstyle=\ttfamily\linespread{\ZSFcodeLeadingFactor}\color{code_fg},
    numbers=none,
    aboveskip=\ZSFspaceM,
    belowskip=\ZSFspaceM,
    frame=single,
    rulecolor=\color{code_rule},
    framerule=\ZSFcodeFrameRule,
    framexleftmargin=\ZSFcodeXMargin,
    framexrightmargin=\ZSFcodeXMargin,
    framextopmargin=\ZSFspaceM,
    framexbottommargin=\ZSFspaceM,
    xleftmargin=\ZSFcodeXMargin,
    xrightmargin=\ZSFcodeXMargin,
    backgroundcolor=\color{code_bg},
    keywordstyle=\bfseries\color{ce_pink},
    commentstyle=\color{ce_green},
    stringstyle=\color{ce_yellow},
    identifierstyle=\color{code_fg},
    showstringspaces=false,
    breaklines=true,
    breakatwhitespace=true,
    literate={->}{{$\rightarrow$}}2 {ü}{{\"u}}1 {ä}{{\"a}}1 {ö}{{\"o}}1 {Ü}{{\"U}}1 {Ä}{{\"A}}1 {Ö}{{\"O}}1 {—}{{---}}1 {−}{{-}}1 {„}{{\glqq}}1 {“}{{\grqq}}1 {”}{{''}}1 {’}{{'}}1 {–}{{--}}1 {→}{{$\rightarrow$}}1
}

% DER STIL GILT AB HIER FÜR JEDES LISTING. Vorher musste ihn jedes einzelne
% mit `[style=CodeExpert]` anfordern — und wer es vergass, bekam kein Fehler,
% keine Warnung und keinen Linter-Treffer, sondern ein Listing ohne Farben,
% ohne Rahmen und mit fremden Abständen. Ein Baustein, dessen Erscheinung
% davon abhängt, ob der Autor eine Option wiederholt hat, ist keine stabile
% API (rules/02_ai_mandate → gleiche Sache, gleiche Bedienung).
% Die Angabe in den Kapiteln bleibt gültig und ist ab jetzt redundant.
\lstset{style=CodeExpert}

% ------------------------------------------------------------
% codebox — Container für Code-Snippets (optional mit Titel)
%   Identisches Erscheinungsbild wie die Titelboxen (zsftitlebox),
%   aber weisser Hintergrund und zero top/bottom padding, da
%   lstlisting eigene above/belowskips mitbringt.
% ------------------------------------------------------------
\tcbset{preset/code/.style={
  zsftitlebox,
  ZSF@surface=plain,
  pady=none,
}}

% ------------------------------------------------------------
% Mehrteilige Code-Gruppen — eine Umgebung mit Regler
% ------------------------------------------------------------
% Eine Box mit vier Kantenformen — also ein Regler, nicht vier Namen
% (rules/02_ai_mandate → Katalog-Ökonomie):
%
%   \begin{codebox}[Titel]                 geschlossene Box (Default)
%   \begin{codebox}[Titel][part=first]     Rahmen oben + Seiten, unten offen
%   \begin{codebox}[][part=mid]            nur Seiten, oben und unten offen
%   \begin{codebox}[][part=last]           Rahmen unten + Seiten, oben offen
%
% Die Teile sind visuell zusammenhängend, physisch getrennt — dazwischen darf
% die Spalte umbrechen. Titel trägt nur der erste Teil.
% ------------------------------------------------------------
% Offene Kanten statt geschlossener Box: eine eigene FORM, aber derselbe
% Vertrag wie jede andere Box. Die Skips sind der bewusste Unterschied — die
% Teile müssen aneinanderstossen.
%
% Genau deshalb `zsfboxcontract` und nicht `zsfboxshape`: Der Vertrag (Ton,
% Justierung, Absatzabstand, Eintritts-Sequenz, Reset der verzögert
% aufgelösten Wahlen) gilt auch für eine Box, die ihre Kanten selbst zeichnet —
% die Form gilt nicht. Wer ihn hier stattdessen abschreibt, lässt früher oder
% später eine Zeile aus, und die Auslassung meldet sich nicht: ein grauer
% Titelbalken statt der Kapitelfarbe sieht für sich plausibel aus.
%
% `ZSF@surface=plain` und `bodyparskip=inherit` sind keine nachgeholten
% Vertragsteile, sondern echte Abweichungen: Code braucht die ungetönte Fläche,
% auf der die Syntaxfarben lesbar bleiben, und kennt keine Absätze.
\tcbset{zsf@codebox@split@base/.style={
  zsfboxcontract,
  unbreakable,
  bodyparskip=inherit,
  ZSF@surface=plain,
  leftrule=\ZSFboxRuleWidth, rightrule=\ZSFboxRuleWidth,
  toprule=0pt, bottomrule=0pt,
  arc=0pt, outer arc=0pt,
  boxsep=0pt,
  padx=normal, pady=none,
  % Der Vorabstand wird über das REGISTER genullt, nicht über `before skip`.
  % Letzteres ist in tcolorbox ein Stil, der `before` schreibt, und überschriebe
  % damit den Eintritts-Hook aus zsfbox (\ZSF@blockBefore): Der Abstand käme an,
  % die Balken-Merker und das \par gingen verloren. Über /utils/exec gilt der
  % Wert nur für diese Box — nachgemessen, der Hook bleibt.
  /utils/exec={\renewcommand{\ZSFBoxBeforeSkip}{0pt}},
  after skip=\ZSFspaceXS,
}}
% Die vier Kantenformen als Auswahl-Schlüssel. 'whole' ist der Default und
% liefert die geschlossene Box; die drei anderen öffnen je eine Kante.
%
% `part` wählt wie `weight` keinen WERT, sondern einen BASISSTIL — und ein
% Basisstil bringt Vorbelegungen mit (`padx=normal`, `pady=none`, Skips).
% Stünde er wie die übrigen Regler in der Aufrufer-Liste, liefen diese
% Vorbelegungen nach den Wahlen des Aufrufers und überschrieben sie:
% `[padx=bar, part=first]` verlor den Innenabstand, `[part=first, padx=bar]`
% behielt ihn — stumm, weil beide Werte für sich gültig sind. Deshalb liest
% die Umgebung die Wahl unten im Vor-Lauf und setzt die fertige Basis VOR die
% Aufrufer-Optionen, genau wie `ZSF@weightBase` in styles/60_boxes.tex.
\tcbset{
  ZSF@partBase/.is choice,
  ZSF@partBase/whole/.style={preset/code},
  ZSF@partBase/first/.style={
    zsf@codebox@split@base,
    toprule=\ZSFboxRuleWidth,
    arc=\ZSFboxInnerArc, outer arc=\ZSFboxArc,
    sharp corners=south,
    % KEIN `before skip` hier. Es ist in tcolorbox ein Stil, der `before`
    % schreibt, und ueberschriebe damit den Eintritts-Hook aus zsfbox
    % (\ZSF@blockBefore) — der Abstand kaeme an, die Balken-Merker gingen
    % verloren. Der Abstand kommt ohnehin von dort, aus demselben Register.
    zsftitlebar,
  },
  ZSF@partBase/mid/.style={zsf@codebox@split@base},
  ZSF@partBase/last/.style={
    zsf@codebox@split@base,
    bottomrule=\ZSFboxRuleWidth,
    arc=\ZSFboxInnerArc, outer arc=\ZSFboxArc,
    sharp corners=north,
    after skip=\ZSFspaceS,
  },
  % Der öffentliche Schlüssel bleibt als No-op bestehen: Er steht weiterhin in
  % der Aufrufer-Liste, wo tcolorbox ihn sonst als unbekannt meldete. Gelesen
  % wurde er zu diesem Zeitpunkt längst.
  part/.is choice,
  part/whole/.code={},
  part/first/.code={},
  part/mid/.code={},
  part/last/.code={},
}
\makeatletter
\newcommand{\ZSF@codeBase}{ZSF@partBase=whole}
\pgfkeys{
  /zsf/codepart/.is family, /zsf/codepart,
  part/.code={\renewcommand{\ZSF@codeBase}{ZSF@partBase=#1}},
  .unknown/.code={},% chktex 26
}
% Öffnet die Box mit EXPANDIERTER Basis — ohne den Umweg landete der Makroname
% selbst in der Optionsliste (dasselbe Muster wie \ZSF@figOpen in 60_boxes).
\newcommand{\ZSF@codeOpen}[3]{\begin{ZSF@box}{#1}{#2}{#3}} % chktex 9
\NewDocumentEnvironment{codebox}{O{} O{}}
  {\renewcommand{\ZSF@codeBase}{ZSF@partBase=whole}%
   \renewcommand{\ZSF@boxOwnKnobs}{,part,}%
   \pgfkeys{/zsf/codepart, #2}%
   \expandafter\ZSF@codeOpen\expandafter{\ZSF@codeBase}{#1}{#2}}
  {\end{ZSF@box}} % chktex 9
\makeatother
